\documentclass[11pt]{article}
\usepackage[utf8]{inputenc}
\usepackage[T1]{fontenc}
\usepackage[spanish,es-nodecimaldot]{babel}
\usepackage{amsmath,booktabs,geometry,graphicx,hyperref,microtype,xcolor}
\geometry{margin=2.35cm}
\hypersetup{colorlinks=true,linkcolor=blue!50!black,urlcolor=blue!60!black}
\title{Forecasting para salidas en RiskGate\\\large Experimento 02: normalización por activo, reducción gradual y boosting cuantílico}
\author{RiskGate Research}
\date{20 de julio de 2026}

\begin{document}
\maketitle
\begin{abstract}
El primer experimento sugirió que GJR--GARCH podía elevar el Sharpe y reducir el drawdown desde 2024, sin
significancia estadística. Esta segunda prueba amplía el test a 2020--2026, normaliza la volatilidad por activo,
reemplaza la salida binaria por estados de exposición 100/50/0 y añade LightGBM cuantílico. El resultado no
confirma la señal inicial: buy-and-hold obtiene Sharpe 0.91, frente a 0.86 para GARCH gradual, 0.80 para LightGBM
y 0.78 para el híbrido. GARCH protege durante 2020--2021, pero empeora 2022 y no conserva la ventaja después.
El intervalo bootstrap al 90\% de la diferencia de Sharpe GARCH--hold es $[-0.45,0.35]$; para LightGBM es
$[-0.24,-0.01]$. Una regla genérica de venta basada sólo en volatilidad o cuantiles diarios no tiene soporte para
integrarse en RiskGate. El siguiente estudio debe condicionar la salida a una operación, horizonte, invalidación
y payoff concretos.
\end{abstract}

\section{Motivación y cambio respecto del experimento 01}
El experimento 01 evaluó diez activos desde 2024. GJR--GARCH--$t$ elevó el Sharpe observado de 1.16 a 1.31,
pero su intervalo bootstrap incluyó ampliamente cero. Además, un umbral común dejó algunos activos siempre
fuera del mercado. La segunda ronda fue diseñada antes de observar sus resultados para atacar esas debilidades:
\begin{enumerate}
  \item comenzar el test en 2020 e incluir varios regímenes;
  \item expresar la volatilidad como percentil de la historia de cada activo;
  \item reducir exposición gradualmente en lugar de liquidar ante un único umbral;
  \item comparar contra regresión cuantílica con gradient boosting;
  \item exigir consistencia por bloques, no sólo una mejora agregada.
\end{enumerate}

\section{Modelos y reglas}
\subsection{GJR--GARCH normalizado}
Se conserva GJR--GARCH(1,1) con innovaciones Student-$t$,
\begin{equation}
 \sigma_t^2=\omega+\alpha\epsilon_{t-1}^2+
 \gamma\epsilon_{t-1}^2\mathbf 1\{\epsilon_{t-1}<0\}+\beta\sigma_{t-1}^2,
\end{equation}
reestimado cada 21 sesiones con un máximo de 1,250 observaciones. Para evitar comparar directamente la
volatilidad de TSLA con la de SPY, $\hat\sigma_{t+1}$ se convierte en un percentil contra las últimas 504
volatilidades realizadas a 20 sesiones del mismo activo. La exposición es
\begin{equation}
a_t^{G}=\begin{cases}
1,&p_t<0.80,\\
0.5,&0.80\le p_t<0.90,\\
0,&p_t\ge0.90.
\end{cases}
\end{equation}
No se optimizan estos cortes en test.

\subsection{LightGBM cuantílico}
Se entrenan tres modelos globales con pérdida pinball para $q\in\{0.1,0.5,0.9\}$. Las covariables incluyen
retornos, volatilidad realizada, momentum, volumen, día de semana, identidad del activo y estado de SPY. La
señal estandarizada es
\begin{equation}
s_t=\frac{\hat q_{0.50,t}}{\max(\hat q_{0.90,t}-\hat q_{0.10,t},10^{-6})}.
\end{equation}
Un único centro de decisión se selecciona en 2019; alrededor de éste se usa una banda fija para asignar
exposición 0, 0.5 o 1. El híbrido toma $\min(a_t^G,a_t^Q)$, es decir, la recomendación más conservadora.

\section{Protocolo}
El universo permanece congelado: SPY, QQQ, IWM, AAPL, MSFT, NVDA, AMZN, GOOGL, META y TSLA. Los precios
ajustados provienen del endpoint público de Yahoo Finance. LightGBM se entrena con 2015--2018, su umbral se
elige solamente en 2019 y el test comienza el 1 de enero de 2020. GARCH usa únicamente historia disponible en
cada fecha. La decisión tomada al cierre $t$ se aplica al retorno $t+1$:
\begin{equation}
r^{strat}_{t+1}=a_t r_{t+1}-0.0005|a_t-a_{t-1}|.
\end{equation}
Se reportan cuatro bloques: 2020--2021, 2022, 2023--2024 y 2025--2026. La incertidumbre de la diferencia de
Sharpe se estima con 2,000 bootstraps de bloques móviles de 20 sesiones.

\section{Resultados agregados}
\begin{table}[ht]
\centering
\caption{Resultados fuera de muestra, 2020--2026.}
\begin{tabular}{lrrrrrr}
\toprule
Estrategia & Ret. anual & Vol. anual & Sharpe & Max DD & Exposición & Cambios\\
\midrule
Buy \& hold & 26.0\% & 28.5\% & 0.91 & $-45.4\%$ & 100.0\% & 0\\
GARCH gradual & 18.7\% & 21.6\% & 0.86 & $-46.4\%$ & 84.5\% & 229\\
LightGBM cuantílico & 22.3\% & 27.8\% & 0.80 & $-46.8\%$ & 96.4\% & 1,108\\
Híbrido & 16.3\% & 21.0\% & 0.78 & $-46.8\%$ & 82.1\% & 1,004\\
\bottomrule
\end{tabular}
\end{table}

\begin{figure}[ht]
\centering
\includegraphics[width=.9\textwidth]{results/experiment_02/equity_curves.pdf}
\caption{Riqueza acumulada fuera de muestra, neta del costo supuesto.}
\end{figure}

Ninguna capa supera buy-and-hold en Sharpe, retorno o drawdown máximo. GARCH disminuye la volatilidad un
24\%, pero sacrifica aproximadamente 7.4 puntos porcentuales de retorno anual y no reduce el peor drawdown.
El mayor número de cambios de LightGBM revela una señal inestable; aun con un costo modesto, la rotación daña
el desempeño. Su intervalo predictivo nominal de 80\% cubre sólo 65.3\%, su mediana acierta el signo 50.4\% y
por tanto está materialmente subcalibrado.

\section{Resultados por régimen}
\begin{table}[ht]
\centering
\caption{Sharpe por bloque.}
\begin{tabular}{lrrrr}
\toprule
Estrategia & 2020--2021 & 2022 & 2023--2024 & 2025--2026\\
\midrule
Buy \& hold & 1.81 & $-1.17$ & 2.84 & 0.67\\
GARCH gradual & 2.24 & $-1.53$ & 2.41 & 0.60\\
LightGBM cuantílico & 1.50 & $-1.28$ & 2.80 & 0.62\\
Híbrido & 1.98 & $-1.63$ & 2.39 & 0.51\\
\bottomrule
\end{tabular}
\end{table}

GARCH funciona en 2020--2021: reduce el drawdown del bloque de $-32.9\%$ a $-14.5\%$ y eleva el Sharpe.
Sin embargo, no protege el mercado bajista persistente de 2022. Una explicación es temporal: un filtro de
volatilidad reacciona después de que comienza el movimiento, reduce exposición tras las pérdidas y puede perder
la recuperación; tampoco distingue volatilidad favorable de deterioro de la tesis. El híbrido hereda las salidas
incorrectas de ambos modelos.

\section{Incertidumbre y decisión experimental}
Los intervalos bootstrap al 90\% de la diferencia de Sharpe contra buy-and-hold son:
\begin{center}
\begin{tabular}{lr}
\toprule
Estrategia & IC 90\%\\
\midrule
GARCH gradual & $[-0.45,\ 0.35]$\\
LightGBM cuantílico & $[-0.24,\ -0.01]$\\
Híbrido & $[-0.58,\ 0.29]$\\
\bottomrule
\end{tabular}
\end{center}
GARCH sigue siendo inconcluso, ahora con mediana negativa. LightGBM es consistentemente peor al nivel usado.
El criterio de éxito no se cumple. No se justifica añadir ninguna de estas reglas al score operativo de RiskGate.

\section{Qué se aprendió}
El resultado no implica que el forecasting carezca de valor. Refuta una formulación concreta: ``alta volatilidad
o cuantil mediano negativo implica vender una posición larga genérica durante la siguiente sesión''. Una decisión
de salida real depende de información omitida:
\begin{itemize}
  \item precio de entrada, beneficio acumulado e invalidación;
  \item horizonte original y tiempo restante;
  \item objetivo y probabilidad de tocar stop antes que objetivo;
  \item estrategia de opción, theta, IV, gamma y payoff;
  \item régimen de mercado y razón por la cual aumentó la volatilidad.
\end{itemize}
El forecasting puede ser más útil para dimensionar riesgo que para decidir dirección. El buen bloque de GARCH
en 2020--2021 sugiere estudiar reducción de tamaño condicionada a una operación, no liquidación universal.

\section{Siguiente experimento}
La siguiente unidad de análisis debe ser una \emph{operación simulada}, no un día de exposición permanente:
\begin{enumerate}
  \item generar entradas con reglas congeladas de RiskGate;
  \item fijar en la entrada stop, objetivo y horizonte máximo;
  \item comparar salida base contra probabilidad de tocar stop antes del objetivo;
  \item usar caminos Monte Carlo condicionados por GARCH y un benchmark histórico no paramétrico;
  \item medir P\&L por operación, hit rate, expected shortfall, tiempo en mercado y falsos exits;
  \item mantener train/valid/test y reportar resultados por régimen y estrategia.
\end{enumerate}
Sólo después de mostrar valor sobre operaciones del tipo que RiskGate realmente evalúa tendría sentido recopilar
históricos de opciones y trasladar la política al payoff multileg.

\section{Conclusión}
La aparente ventaja inicial de GJR--GARCH no fue robusta al ampliar el horizonte y normalizar por activo.
LightGBM cuantílico tampoco añadió valor y mostró mala calibración. Este resultado evita integrar prematuramente
una regla que habría reducido retorno sin controlar el peor drawdown. La línea de investigación continúa, pero
la pregunta correcta cambia de ``¿cuándo salir del mercado?'' a ``dada esta operación concreta, ¿cambió la
probabilidad de alcanzar su invalidación antes que su objetivo?''.
\end{document}
